\section{题06：独立三角形（P16788）}

\subsection{问题重述}

二维平面上有 $N$ 个互不重合的整数点，每个点的纵坐标只可能是 $1$ 或 $-1$。我们需要用这些点连出尽可能多的三角形，每个三角形须满足三个顶点不共线，且不同三角形之间不能共享顶点——一个点一旦作为三角形的顶点，就已经连接到该三角形内部的另外两个顶点，共消耗了两条线段的端点，这也意味着每个点至多参与一个三角形（因为每个点最多作为两条线段的端点）。在三角形数量取到最大的所有方案中，进一步要求所有三角形的面积之和最大。输入保证该最大面积和为整数。$N \le 2 \times 10^5$，$|x_i| \le 10^9$。

\subsection{建模：当所有点只分布在两条水平线上}

这是一个典型的“受限组合优化”问题。输入本质上是两组一维坐标的集合：上方点集（$y = 1$）的横坐标列表和下方点集（$y = -1$）的横坐标列表。设上方点个数为 $A$，下方点个数为 $B$，则 $A + B = N$。

由于所有点只能位于 $y = 1$ 和 $y = -1$ 两条直线上，而三个顶点若共线则不能构成三角形，因此任何合法三角形的三个顶点不可能全在上方或全在下方。每个三角形的顶点构成只可能有两种情形：两个上方点搭配一个下方点（以下称为“以底在上”的三角形），或一个上方点搭配两个下方点（“以底在下”的三角形）。两种情形中，三角形的面积公式都是 $\frac12 \times \text{底边长度} \times 2 = \text{底边长度}$，因为底边由同侧的两个顶点确定，其对侧顶点到底边的垂直距离恒为 $2$。换言之，三角形的面积恰好等于其底边两端点的横向距离，与第三条顶点的具体位置无关。

\subsection{朴素思路：暴力枚举为何不可行}

如果采用最暴力的做法，可以从 $N$ 个点中枚举所有可能的三元组，检查其是否构成合法三角形，再在合法三角形集合中找出若干个两两不共享顶点的方案。枚举所有三角形的时间至少是 $\mathrm{O}(N^3)$，在 $N = 2 \times 10^5$ 的规模下完全不可行。即便我们利用“所有点仅在两条线上”这一性质把候选三角形压缩到“二上一下”和“一上二下”两类，每类的候选数量依然是 $\mathrm{O}(N^2)$ 级别——上方点的每对组合都可以搭配任意一个下方点。因此必须寻找更深刻的结构性结论，将指数级的组合选择转化为简单的线性规划与贪心。

\subsection{核心洞察：每个三角形消耗一上一下两条边}

设我们最终构造了 $x$ 个以底在上的三角形和 $y$ 个以底在下的三角形，总三角形数 $T = x + y$。从顶点消耗的角度看：每个以底在上的三角形消耗 $2$ 个上方点和 $1$ 个下方点，每个以底在下的三角形消耗 $1$ 个上方点和 $2$ 个下方点。因此上方点的总消耗为 $2x + y$，下方点的总消耗为 $x + 2y$，它们分别不能超过 $A$ 和 $B$：

\begin{equation}\label{eq:constraints}
2x + y \le A, \qquad x + 2y \le B, \qquad x \ge 0,\; y \ge 0.
\end{equation}

将两式相加得到 $3(x + y) \le A + B = N$，即 $T = x + y \le \lfloor N/3 \rfloor$。另一方面，每个三角形无论哪种类型都至少需要 $1$ 个上方点和 $1$ 个下方点，所以 $T \le \min(A, B)$。综合这两条独立的上界，我们得到三角形数量的理论上界：

\begin{equation}\label{eq:maxT}
T_{\max} = \min\!\big(\lfloor N/3 \rfloor,\; A,\; B\big).
\end{equation}

这个上界有没有意义？以样例 $1$ 为例，$N = 3$，三个点分别为 $(0, 1), (2, 1), (1, -1)$，则 $A = 2$，$B = 1$。$T_{\max} = \min(1, 2, 1) = 1$，确实可以达到 $1$ 个三角形。样例 $2$ 中 $N = 6$，上方点坐标 $0, 1, 3$（$A = 3$），下方点坐标 $1, 2, 4$（$B = 3$），$T_{\max} = \min(2, 3, 3) = 2$，同样可达到。

\subsection{上界的可达性证明：线性不等式组的整数解}

上界给出的是必要条件，我们还需证明它总是充分的——即对于任意满足 $T = T_{\max}$ 的 $T$，总存在非负整数 $x, y$ 使得 $x + y = T$ 且满足不等式组~(\ref{eq:constraints})。将 $y = T - x$ 代入，约束转换为 $x$ 的一个闭区间：

\[
2x + (T - x) \le A \;\Longrightarrow\; x \le A - T,
\]
\[
x + 2(T - x) \le B \;\Longrightarrow\; x \ge 2T - B.
\]

再加上 $0 \le x \le T$，得到 $x$ 的可行区间：

\begin{equation}\label{eq:xrange}
\max(0,\, 2T - B) \le x \le \min(T,\, A - T).
\end{equation}

该区间非空等价于左右端点满足 $\max(0,\, 2T - B) \le \min(T,\, A - T)$。分情况验证：若 $T = \lfloor N/3 \rfloor$ 且 $T \le A, B$，由 $3T \le A + B$ 可推出 $2T - B \le A - T$，同时 $A - T \ge 0$（因为 $T \le A$），且 $2T - B \le T$（因为 $T \le B$），故区间非空。若 $T = A$（此时 $A$ 是三者中最小者），代入得 $T = A \le \lfloor N/3 \rfloor$，结合 $A + B = N$ 可知 $2A \le B$，此时可行区间退化为 $\max(0, 2A - B) \le x \le \min(A, 0)$，即 $x = 0$——全部三角形均为以底在下的类型。对称地，若 $T = B$ 则 $x = T$——全部为以底在上的类型。由此可见，$T_{\max}$ 总是可以通过合理分配两种三角形的数量来实现。

\subsection{最大化面积之和：底边配对的最优策略}

由于每个三角形的面积等于其底边两端点的横向距离，而第三条顶点的横坐标不影响面积，因此最大化面积之和的问题可以拆解为：从 $A$ 个上方点中选出 $2x$ 个点，两两配成 $x$ 条底边，使这 $x$ 条底边长度之和最大；同时对下方点做同样操作（选出 $2y$ 个点配成 $y$ 条底边）。剩下的 $A - 2x$ 个上方点和 $B - 2y$ 个下方点充当第三条顶点，它们的横坐标无关紧要。

现在问题简化为：给定 $m$ 个已排序的点（横坐标从小到大），要从中选出 $2k$ 个点配成 $k$ 对，使得每对的横向距离之和最大。直觉上，要使距离最大，应该让选出的点尽可能分散——最左和最右配对、次左和次右配对，依此类推。这一直觉可以用重排不等式严格证明：对于已排序数组 $p_0 \le p_1 \le \dots \le p_{m-1}$，在所有选出 $2k$ 个元素并配对的方案中，配对方式 $(p_0, p_{m-1}), (p_1, p_{m-2}), \dots, (p_{k-1}, p_{m-k})$ 使得 $\sum (p_{m-i} - p_i)$ 取到最大值。原因在于：内层的点无论如何选，其贡献的距离不会超过外层点产生的跨度；如果我们从两端选取最远的 $k$ 对，中间的点即使被选为第三条顶点，也不会影响底边长度，而任何试图使用中间点作为底边端点的方案都会使距离和变小。

以样例 $2$ 的上方点集 $(0, 1, 3)$ 为例：$A = 3$，若 $x = 1$，需选出 $2$ 个点配对。选取最左的 $0$ 和最右的 $3$，底边长度为 $3$；若选 $0, 1$ 或 $1, 3$，长度仅为 $1$ 或 $2$，均不如 $3$。

基于这一结论，我们分别将上方点和下方点按横坐标升序排列后，预计算“前缀最大宽度和”：

\[
\mathit{topW}[k] = \sum_{i=0}^{k-1} (t_{A-1-i} - t_i),\qquad
\mathit{botW}[k] = \sum_{i=0}^{k-1} (b_{B-1-i} - b_i).
\]

其中 $\mathit{topW}[k]$ 表示从上方点中选出 $2k$ 个并以最优方式配成 $k$ 条底边所能获得的最大底边长度之和，$\mathit{botW}[k]$ 含义对称。注意 $\mathit{topW}[k]$ 仅在 $2k \le A$ 时有定义（同理 $\mathit{botW}[k]$ 要求 $2k \le B$）；超出此范围的 $k$ 取值为 $0$（无法构造那么多以该侧为底的三角形）。

有了这两个前缀和数组，对于一个可行的 $(x, y)$ 组合（其中 $y = T - x$，$x$ 满足区间~(\ref{eq:xrange})），该方案的总面积即为 $\mathit{topW}[x] + \mathit{botW}[y]$。遍历 $x$ 的所有可行取值并记录最大值，即可得到在 $T = T_{\max}$ 个三角形的前提下所能获得的最大面积之和。

\subsection{复杂度分析}

整个算法的瓶颈在于排序：分别对上方点集和下方点集按横坐标排序，时间复杂度为 $\mathrm{O}(A \log A + B \log B) = \mathrm{O}(N \log N)$。前缀和数组的计算是线性的 $\mathrm{O}(\max T)$，遍历 $x$ 的可行区间同样是 $\mathrm{O}(\max T) \le \mathrm{O}(N)$。在 $N \le 2 \times 10^5$ 的数据范围下，$\mathrm{O}(N \log N)$ 的排序足以在时限内完成。空间复杂度为 $\mathrm{O}(N)$，用于存储两个坐标数组和两个前缀和数组。

\subsection{实现细节}

代码实现中有几个值得注意的细节。首先，由于 $|x_i| \le 10^9$ 且 $N$ 可达 $2 \times 10^5$，面积之和可能超过 $32$ 位整型的表示范围，因此前缀和数组 $\mathit{topW}$ 和 $\mathit{botW}$ 以及最终答案 $\mathit{bestArea}$ 必须使用 \mil{long long}（$64$ 位有符号整数）。其次，在计算前缀和时，需要确保索引不越界：\mil{topW[k]} 仅当 $2k \le A$ 且 $k \le \max T$ 时才可计算；代码中通过 \mil{min(A/2, maxT)} 来控制循环上界，对于超出此范围的 $k$，$\mathit{topW}[k]$ 和 $\mathit{botW}[k]$ 保持初始化时的 $0$，这恰好对应“无法构造 $k$ 个以该侧为底的三角形则贡献为 $0$”的语义。此外，$T_{\max}$ 可能为 $0$（例如所有点都在同一条线上），此时需特判并输出 \mil{0 0}，避免后续计算中出现空的循环范围导致未定义行为。最后，$x$ 的可行区间由 $T_{\max}$ 和 $A, B$ 共同确定：下界 $\max(0, 2T - B)$ 保证下方点够用，上界 $\min(T, A - T)$ 保证上方点够用；这两条边界均来源于不等式组~(\ref{eq:constraints})的直接推导。

\subsection{代码实现}
\inputminted{c++}{../src/p06_independent_triangles.cpp}
